% =========================================================
% Interval Arithmetic in $\mathbb{R}^n$
% =========================================================

\section{\texorpdfstring{Interval Arithmetic in $\mathbb{R}^n$}{Interval Arithmetic in R n}}
\label{sec:interval-arithmetic-rn}

\begin{definitionbox}{Definition (Box / Product Interval)}
\begin{definition}[Box / Product Interval]\label{def:box}
Let $a,b\in\mathbb{R}^n$ with $a\le b$ coordinatewise (i.e., $a_i\le b_i$ for all $i$).
The \emph{box} (or \emph{interval vector}) determined by $(a,b)$ is
\[
[a,b] \;:=\; \{x\in\mathbb{R}^n : a\le x \le b\}
\;=\;\prod_{i=1}^n [a_i,b_i].
\]
\end{definition}
\end{definitionbox}

\begin{dependencies}
\begin{itemize}
  \item \hyperref[def:cartesian-product]{Cartesian Product}
  \item \hyperref[def:closed-interval]{Closed Interval}
  \item \hyperref[def:order-on-r]{Order on $\mathbb{R}$}
\end{itemize}
\end{dependencies}

\begin{definitionbox}{Definition (Box Containment)}
\begin{definition}[Box Containment]\label{def:box-contain}
For $a\le b$ and $c\le d$ in $\mathbb{R}^n$, we write $[a,b]\subseteq[c,d]$ iff
\[
(\forall x\in\mathbb{R}^n)\ \bigl(x\in[a,b]\Rightarrow x\in[c,d]\bigr).
\]
Equivalently (and most useful computationally),
\[
[a,b]\subseteq[c,d]\quad \Longleftrightarrow\quad c\le a\ \land\ b\le d
\qquad\text{(coordinatewise).}
\]
\end{definition}
\end{definitionbox}

\begin{dependencies}
\begin{itemize}
  \item \hyperref[def:box]{Box / Product Interval}
  \item \hyperref[def:subset]{Subset}
  \item \hyperref[def:order-on-r]{Order on $\mathbb{R}$}
\end{itemize}
\end{dependencies}

\begin{definitionbox}{Definition (Minkowski Addition of Boxes)}
\begin{definition}[Minkowski Addition of Boxes]\label{def:box-add}
For boxes $X=[a,b]$ and $Y=[c,d]$ in $\mathbb{R}^n$, their \emph{Minkowski sum} is
\[
X\oplus Y \;:=\; \{x+y : x\in X,\ y\in Y\}.
\]
For boxes, this closes in the class of boxes and satisfies the endpoint rule
\[
[a,b]\oplus[c,d] \;=\; [a+c,\; b+d].
\]
\end{definition}
\end{definitionbox}

\begin{dependencies}
\begin{itemize}
  \item \hyperref[def:box]{Box / Product Interval}
  \item \hyperref[def:addition-on-r]{Addition on $\mathbb{R}$}
\end{itemize}
\end{dependencies}

\begin{definitionbox}{Definition (Box Subtraction / Difference)}
\begin{definition}[Box Subtraction / Difference]\label{def:box-sub}
For boxes $X=[a,b]$ and $Y=[c,d]$ in $\mathbb{R}^n$, define
\[
X\ominus Y \;:=\; \{x-y : x\in X,\ y\in Y\}.
\]
Then
\[
[a,b]\ominus[c,d] \;=\; [a-d,\; b-c].
\]
\end{definition}
\end{definitionbox}

\begin{dependencies}
\begin{itemize}
  \item \hyperref[def:box]{Box / Product Interval}
  \item \hyperref[def:subtraction-on-r]{Subtraction on $\mathbb{R}$}
\end{itemize}
\end{dependencies}

\begin{definitionbox}{Definition (Scalar Multiplication of a Box)}
\begin{definition}[Scalar Multiplication of a Box]\label{def:box-smul}
For $\lambda\in\mathbb{R}$ and a box $[a,b]\subseteq\mathbb{R}^n$, define
\[
\lambda[a,b] \;:=\; \{\lambda x : x\in[a,b]\}.
\]
Then
\[
\lambda[a,b]=
\begin{cases}
[\lambda a,\ \lambda b], & \lambda\ge 0,\\[4pt]
[\lambda b,\ \lambda a], & \lambda<0,
\end{cases}
\]
where the inequalities are coordinatewise.
\end{definition}
\end{definitionbox}

\begin{dependencies}
\begin{itemize}
  \item \hyperref[def:box]{Box / Product Interval}
  \item \hyperref[def:multiplication-on-r]{Multiplication on $\mathbb{R}$}
  \item \hyperref[def:order-on-r]{Order on $\mathbb{R}$}
\end{itemize}
\end{dependencies}

\begin{definitionbox}{Definition (Coordinatewise / Hadamard Product of Boxes)}
\begin{definition}[Coordinatewise / Hadamard Product of Boxes]\label{def:box-hadamard}
For $x,y\in\mathbb{R}^n$, write $(x\odot y)_i := x_i y_i$.
For boxes $X=[a,b]$ and $Y=[c,d]$, define
\[
X\odot Y \;:=\; \{x\odot y : x\in X,\ y\in Y\}.
\]
Then $X\odot Y$ is again a box, computed coordinatewise:
\[
[a,b]\odot[c,d] \;=\; \prod_{i=1}^n \Bigl[\min S_i,\ \max S_i\Bigr],
\qquad
S_i:=\{a_ic_i,\ a_id_i,\ b_ic_i,\ b_id_i\}.
\]
\end{definition}
\end{definitionbox}

\begin{dependencies}
\begin{itemize}
  \item \hyperref[def:box]{Box / Product Interval}
  \item \hyperref[def:multiplication-on-r]{Multiplication on $\mathbb{R}$}
  \item \hyperref[def:order-on-r]{Order on $\mathbb{R}$}
\end{itemize}
\end{dependencies}

\begin{remark*}[Interpretation: intervals as guaranteed enclosures]
A box $X=[a,b]\subseteq\mathbb{R}^n$ represents a \emph{set of possible vectors}
(e.g., a computed vector plus uncertainty bounds). The operations above are
\emph{set operations} (Minkowski sum/difference, scalar image), so the result is a
\emph{guaranteed enclosure} of all possible true outcomes under the corresponding
vector operation.
\end{remark*}

\begin{remark*}[Norm bounds and why boxes are convenient]
For many applications you also want quick bounds on a norm over a box.
For example, if $0\le a\le b$ coordinatewise, then for every $x\in[a,b]$ we have
$0\le x\le b$ and hence $\|x\|_2\le \|b\|_2$. More generally, bounding norms over
a box can be reduced to checking finitely many ``corners'' when the objective is
coordinatewise monotone.
\label{rem:norm-bounds}
\end{remark*}

\begin{remark*}[Status]
\textit{This section is a stub. A full treatment -- interval extensions of general
functions $f:\mathbb{R}^n\to\mathbb{R}^m$, the dependency problem, sub-distributivity,
natural vs. centered (affine) forms, and connections to compactness and continuity --
will be developed in a future revision.}
\end{remark*}
